\documentclass{article}
\usepackage[utf8]{inputenc}
\usepackage[T1]{fontenc}
\usepackage{geometry}
\usepackage{graphicx}
\usepackage{float}
\usepackage{hyperref}
\usepackage[english,indonesian]{babel}

\title{Optimasi Pemrosesan Template LaTeX Menggunakan Pendekatan Berbasis AI}
\author{Academic Support Agent}
\date{\today}

\begin{document}
\maketitle

\begin{abstract}

\end{abstract}

\chapter{PENDAHULUAN}


\chapter{TINJAUAN PUSTAKA}


% Original content (commented out):
% \documentclass[dvipsnames,a4paper,12pt]{article}
\usepackage[brazil]{babel}
\usepackage[T1]{fontenc}
\usepackage[utf8x]{inputenc}
\usepackage{graphicx}
\usepackage{hyperref}
\usepackage{color}
\usepackage{amsmath}
\usepackage{listings} 
\usepackage[section]{placeins}

\begin{document}
\begin{titlepage}
\begin{center}
   \includegraphics[width=4cm]{logoicexpreto}\par
	\vspace{1cm}
	{\scshape\Large\bfseries Laboratório de Física Computacional I \par}
	\vspace{2.5cm}
    {\huge\bfseries Título \par}
	\vspace{.5cm}
	{\LARGE\itshape Al Uno \par}
    \vspace{.5cm}    
    {\large\itshape Matrícula \par}
    \vspace{.5cm}    
    {\Large\itshape Curso de Física}
    \vspace{1.5cm}    
    \begin{abstract}
    Resumo (3 linhas)
    \end{abstract}
    \vfill
   {\large \today\par}
\end{center}
\end{titlepage}

\tableofcontents{}

\section{Introdução}\label{sec:intro}
Escrito em \LaTeX~\cite{lamport94}. Resultados apresentados na seção~\ref{sec:results}.

\section{Problema}
\begin{figure}[ht] 
        \includegraphics[width=0.5\textwidth]{logoicexpreto}
        \caption{
                \label{fig:logoicex} 
                \textbf{ Toda figura tem legenda. Todos os eixos devem ter rótulos.}
        }
\end{figure}


\section{Resultados}\label{sec:results}
\newpage

\section*{Apêndice}

\definecolor{mygreen}{rgb}{0,0.6,0}
\definecolor{mygray}{rgb}{0.5,0.5,0.5}
\definecolor{mymauve}{rgb}{0.58,0,0.82}

\lstset{ %
  backgroundcolor=\color{white},   % choose the background color; you must add \usepackage{color} or \usepackage{xcolor}
  basicstyle=\small,        % the size of the fonts that are used for the code
  breakatwhitespace=false,         % sets if automatic breaks should only happen at whitespace
  breaklines=true,                 % sets automatic line breaking
  captionpos=b,                    % sets the caption-position to bottom
  commentstyle=\color{mygreen},    % comment style
  escapeinside={\%*}{*)},          % if you want to add LaTeX within your code
  extendedchars=true,              % lets you use non-ASCII characters; for 8-bits encodings only, does not work with UTF-8
  keepspaces=true,                 % keeps spaces in text, useful for keeping indentation of code (possibly needs columns=flexible)
  keywordstyle=\color{blue},       % keyword style
  language=C,                 	   % the language of the code
  numbers=left,                    % where to put the line-numbers; possible values are (none, left, right)
  numbersep=5pt,                   % how far the line-numbers are from the code
  numberstyle=\tiny\color{mygray}, % the style that is used for the line-numbers
  rulecolor=\color{black},         % if not set, the frame-color may be changed on line-breaks within not-black text (e.g. comments (green here))
  showspaces=false,                % show spaces everywhere adding particular underscores; it overrides 'showstringspaces'
  showstringspaces=false,          % underline spaces within strings only
  showtabs=false,                  % show tabs within strings adding particular underscores
%  stepnumber=2,                    % the step between two line-numbers. If it's 1, each line will be numbered
  stringstyle=\color{mymauve},     % string literal style
  tabsize=2,	                   % sets default tabsize to 2 spaces
  title=\lstname                   % show the filename of files included with \lstinputlisting; also try caption instead of title
}


\begin{lstlisting}
#include <stdio.h>
#define N 10
/* Block
 * comment */

int main()
{
    int i;

    // Line comment.
    puts("Hello world!"); 
    for (i = 0; i < N; i++)
    {
        puts("LaTeX is also great for programmers!");
    }
    return 0;
}
\end{lstlisting}

\bibliographystyle{apalike}
\begin{thebibliography}{}
\bibitem[Lamp94]{lamport94} Leslie Lamport, \emph{\LaTeX: a document preparation system}, Addison Wesley, Massachusetts, 2nd edition, 1994.
\end{thebibliography}
\end{document}

\end{document}